\documentclass[12pt,a4paper]{article}
\usepackage[margin=22mm]{geometry}
\usepackage[T1]{fontenc}
\usepackage[utf8]{inputenc}
\setlength{\parindent}{0pt}
\setlength{\parskip}{0.8em}
\linespread{1.15}\selectfont

\begin{document}

\begin{center}
{\LARGE Presentation 2 Script}
\end{center}

Hello everyone.

My name is Heizo Nakai, and my presenter number is 55.

Today, I'd like to introduce my research theme: ``Visual Noise Cancelling VR device.''

Let me start with a simple question.
Could you concentrate on your work at a cluttered desk like this?

I think many people would feel distracted.
I feel the same way.

Then, how about this situation?
Here, the visual noise has been removed.
I think this looks like a much better environment for concentration.

My goal is to enable users to remove visual noise through a VR device.

Next, I will explain the detailed process for visual noise cancellation.

The process consists of two main steps.

The first step is object removal.
The user selects the object they want to remove,
 and the system removes it from the scene.
The second step is inpainting the removed area using AI.
This is similar to Google's ``Magic Eraser'' function.

I aim to implement this kind of processing on a VR device.
For a natural user experience, it must run in real time.

However, VR devices have limited computational resources and battery capacity.

Therefore, the system requires high-speed processing and low power consumption
under these constraints.

Finally, I will explain my approach to addressing these challenges.

I am considering designing a dedicated accelerator.

First, I will analyze the computational cost of each processing step.
Then, I will identify the bottleneck, which is the most time-consuming part of the pipeline.

Based on this analysis, I will design a dedicated accelerator for the bottleneck process.
Through this approach, I aim to reduce both processing time and power consumption.

Thank you for listening.

\end{document}
